problem:  
Let \( (X,\omega,\Omega) \) be a compact Calabi–Yau threefold, and suppose \( \pi: X \to B \) is a special Lagrangian torus fibration outside a singular set \( \Sigma \subset B \) of real codimension two, with smooth torus fibers \( L_b = \pi^{-1}(b) \) for \( b \in B \setminus \Sigma \).  
On \( B \setminus \Sigma \), the McLean theory provides a natural affine structure and a Monge–Ampère metric \( g_B = \mathrm{Hess}(\phi) \) arising from period integrals of \( \Omega \).  

Near the singular locus \( \Sigma \), special Lagrangian limits and disk instanton corrections are expected to modify both the affine structure and its associated potential \( \phi \), producing the **quantum-corrected metric** anticipated by mirror symmetry.

**Problem:**  
Assume \( \Sigma \) is a smooth real surface sitting in \(B\), and that in a neighborhood \( U \subset B\) of a generic point of \( \Sigma \), the special Lagrangian fibers of \( \pi^{-1}(U\setminus\Sigma) \) collapse with locally modeled singularities (for instance of type \(T^2 \times \text{cone}(S^1)\)).   

Is it possible, purely from the **geometric analysis on \(X\)** and the **asymptotics of collapsing special Lagrangian metrics**, to define a canonical “corrected potential” \( \phi_{\mathrm{corr}} \), a real analytic function on \( U \setminus \Sigma \) satisfying  
\[
\det(\mathrm{Hess}\,\phi_{\mathrm{corr}}) = 1 + \sum_{\beta} n_{\beta} e^{-\mathrm{Area}(\beta)},
\]
where the correction coefficients \(n_\beta\) encode counts of holomorphic disks bounded by singular special Lagrangian fibers approaching a class \(\beta\)?   
Equivalently, can one describe an *analytic mechanism* (independent of algebro-geometric models) through which corrections to the semi-flat Monge–Ampère structure naturally arise from the degeneration behavior of the special Lagrangian foliation near \(\Sigma\)?

Why is it a "good" problem:  
This reformulation targets the **core analytic riddle** at the heart of the Strominger–Yau–Zaslow program: how instanton corrections and singular fiber degenerations modify the pure differential-geometric structure of a special Lagrangian base. Unlike the classical semi-flat regime, where everything is smooth and the duality is well understood, this problem directly addresses the real geometric source of mirror symmetry’s quantum corrections within a rigorous PDE-analytic setting.  

The question seeks a purely differential-geometric explanation for the appearance of holomorphic disk contributions in the corrected Monge–Ampère equation. A resolution would forge a direct analytic bridge between calibrated geometry, minimal submanifold theory, and enumerative geometry, establishing the first geometric derivation of mirror “instanton corrections” from degeneration limits.  

Its formulation is concise—merely asking for a corrected potential determined by degenerating special Lagrangian fibrations—but its depth lies in reconciling geometric analysis (collapsing metrics, singular PDEs) with symplectic enumeration. Such a problem cuts to the conceptual core of how mirror symmetry emerges from real Calabi–Yau geometry, making it a quintessential "good" and genuinely open research-level problem.